\section{What is in this repository?}

This repository is a PyTorch research codebase for \textbf{video watch-time prediction}. It contains an implementation of the paper:

\begin{quote}
\textit{``Multi-Granularity Distribution Modeling for Video Watch Time Prediction via Exponential-Gaussian Mixture Network''}
\end{quote}

The repository also includes several baseline models and code to preprocess the public \href{https://kuairec.com/}{KuaiRec} dataset.

\subsection{High-level structure}
\begin{itemize}[leftmargin=*]
  \item \lstinline|model/|: Model implementations (proposed EGMN and baselines).
  \item \lstinline|dataloader/|: Dataset loaders that read preprocessed files and provide PyTorch \lstinline|DataLoader|s.
  \item \lstinline|dataset/kuairec/|: KuaiRec preprocessing scripts.
  \item \lstinline|run_*.py|: Entrypoints to train/evaluate each method.
  \item \lstinline|utils.py|: Evaluation metrics (MAE, XAUC, KL divergence) and helper utilities.
\end{itemize}

\subsection{Code examples}
\paragraph{Typical entrypoint invocation}
\begin{lstlisting}[language=bash]
# Train + evaluate the proposed EGMN model on KuaiRec
python run_egmn.py --dataset_name kuairec --device cuda:0

# Baselines (examples)
python run_cread.py --dataset_name kuairec
python run_d2q.py --dataset_name kuairec
\end{lstlisting}

\paragraph{What the dataloader yields}
The KuaiRec dataloader returns a tuple \lstinline|(features_dict, label)|, where \lstinline|features_dict| contains tensors keyed by feature name (e.g. \lstinline|duration|, categorical IDs, and sequence masks).
\begin{lstlisting}[language=Python]
# Pseudocode reflecting dataloader/kuairec.py
for features, label in dataloaders["train"]:
    # features is a dict of tensors, label is the normalized play_time
    y_hat = model.predict(features)
\end{lstlisting}

\paragraph{Metrics used in evaluation}
\begin{lstlisting}[language=Python]
# utils.py (simplified)
def eval_mae(labels, scores):
    return np.mean(np.abs(labels - scores))

def eval_xauc(labels, preds):
    # ranks by prediction, then measures label ordering agreement
    ...

def eval_kl(samples_p, samples_q, bins=100, epsilon=1e-10):
    hist_p, bin_edges = np.histogram(samples_p, bins=bins, density=True)
    hist_q, _ = np.histogram(samples_q, bins=bin_edges, density=True)
    ...
    return np.sum(hist_p * np.log(hist_p / hist_q))
\end{lstlisting}

\subsection{Evaluation outputs}
The run scripts train on \lstinline|train| and evaluate on \lstinline|test|, printing (at least):
\begin{itemize}[leftmargin=*]
  \item \textbf{MAE}: mean absolute error on watch-time prediction.
  \item \textbf{XAUC}: a ranking-style metric computed from label--prediction ordering.
  \item \textbf{KL}: KL divergence between binned empirical distributions.
\end{itemize}
